\section{Several quadratic factors: the complete second derivative and confluence}
\label{mm:section}

The seventh further-research question in \cite{QSpec} asks for the
second derivative of a spectrum with several masses and its coincident
limit.  The even-polynomial setting of that report admits a complete
answer.  In contrast to the first derivative, a logarithmic cross term
remains; every coefficient of its coincident-mass expansion nevertheless
has a finite polygamma and Bernoulli evaluation.

\subsection{A finite definition for the weighted logarithmic kernel}

Fix $a>0$ and $P(v)=\sum_{d=0}^{D}p_dv^d$.  For $B,C>0$ set
\begin{equation}\label{mm:eq:beta}
 \beta_k(B,C)=(-1)^k\sum_{j=1}^{k-1}
        \frac{(B^j-C^j)(B^{k-j}-C^{k-j})}{j(k-j)},\qquad k\ge2.
\end{equation}
These are the coefficients in
\[
 \log^2\frac{1+Bz}{1+Cz}=\sum_{k\ge2}\beta_k(B,C)z^k.
\]
They are symmetric in $B,C$ and vanish when $B=C$.  Define the
following specific compensated kernel, with $x_n=n+a$:
\begin{align}
 K_{P,a}(B,C)
 &=\sum_{n\ge0}P(x_n^2)
    \left\{\log^2\frac{x_n^2+B}{x_n^2+C}
                 -\sum_{k=2}^{D}\frac{\beta_k(B,C)}{x_n^{2k}}\right\}
    \notag\\
 &\quad+\sum_{d=0}^{D}\sum_{k=2}^{D}
              p_d\beta_k(B,C)\zeta(2k-2d,a).
              \label{mm:eq:K}
\end{align}
All inner sums are empty when $D<2$.  The infinite sum has
$O(x_n^{-2})$ tails and is ordinary and locally normally convergent.
The finite counterterms have no pole: their Hurwitz orders are even,
hence never one.  For $P=1$ or $P(v)=v$, the unmodified logarithmic
sum already converges absolutely.

\begin{theorem}[All even weights and any number of masses]
\label{mm:thm:second}
For $q\ge1$ and $\mathbf B=(B_1,\ldots,B_q)$ with each $B_j>0$,
let
\begin{align}
 \mathcal Z_{P,\mathbf B}(s;a)
 &=\sum_{n\ge0}P(x_n^2)\prod_{j=1}^{q}(x_n^2+B_j)^{-s},
 \label{mm:eq:product}\\
 Z_{P,j}(s;a)&=\sum_{n\ge0}P(x_n^2)(x_n^2+B_j)^{-s}.
 \label{mm:eq:individual}
\end{align}
Their continuations are holomorphic at zero, and
\begin{equation}\label{mm:eq:second}
 \boxed{\displaystyle
 \mathcal Z_{P,\mathbf B}''(0;a)
 =q\sum_{j=1}^{q}Z_{P,j}''(0;a)
             -\sum_{1\le i<j\le q}K_{P,a}(B_i,B_j).}
\end{equation}
Repeated masses are allowed.  Each coincident pair contributes zero.
\end{theorem}
\begin{proof}
Compare the product spectrum with
$q^{-1}\sum_jZ_{P,j}(qs;a)$.  At each lattice point write
$L_j=\log(x^2+B_j)$.  The constant and first-order coefficients of
\[
 \frac1q\sum_j e^{-qsL_j}-e^{-s\sum_jL_j}
\]
vanish.  Its second derivative at zero is the elementary variance
identity
\begin{equation}\label{mm:eq:variance}
 q\sum_jL_j^2-\left(\sum_jL_j\right)^2
                =\sum_{i<j}(L_i-L_j)^2.
\end{equation}
To justify this calculation after continuation, factor out $x^{-2qs}$
and expand the remaining factors at infinity.  Only even powers of
$x^{-1}$ occur, including after multiplication by $P(x^2)$.  A pole
at $s=0$ would require a resulting $x^{-1}$ term; there is none.
Subtract finitely many asymptotic terms to leave a normally summable
remainder.  The constant and linear coefficients of each difference
term vanish before the Hurwitz factors are differentiated.  Hence its
second derivative is exactly the finite-subtraction continuation of
\eqref{mm:eq:variance}.  The latter is
\eqref{mm:eq:K} for each pair, proving \eqref{mm:eq:second}.
The same finite-subtraction argument proves holomorphy at zero for
the original spectra.  Repeated parameters follow either directly
or from local normal convergence.
\end{proof}

For two unweighted factors, this reads
\[
 \mathcal Z_{1,(B,C)}''(0;a)
 =2Z_{1,B}''(0;a)+2Z_{1,C}''(0;a)
       -\sum_{n\ge0}\log^2\frac{x_n^2+B}{x_n^2+C}.
\]
At $B=C$ it reduces to $4Z_{1,B}''(0;a)$, exactly as required by
$\mathcal Z_{1,(B,B)}(s;a)=Z_{1,B}(2s;a)$.

\subsection{An explicit digamma kernel determines the whole correction}

The elementary quadratic resolvent is
\begin{equation}\label{mm:eq:R}
 R_a(B)=\sum_{n\ge0}\frac1{x_n^2+B}
       =\frac{\psi(a+i\sqrt B)-\psi(a-i\sqrt B)}{2i\sqrt B}.
\end{equation}
For a polynomial weight define its finite continuation
\begin{equation}\label{mm:eq:RP}
 R_{P,a}(B)=P(-B)R_a(B)
    +\sum_{j=1}^{D}p_j\sum_{\ell=0}^{j-1}
                      (-B)^\ell\zeta(-2(j-1-\ell),a).
\end{equation}
It is the continued value $Z_{P,B}(1;a)$.  Formula
\eqref{mm:eq:RP} follows from polynomial division of
$x^{2j}/(x^2+B)$ and the Hurwitz values at nonpositive even integers.
Those values are Bernoulli polynomials:
$\zeta(-m,a)=-B_{m+1}(a)/(m+1)$ \cite{DLMFZeta}.

\begin{theorem}[A normalized double primitive of a digamma divided difference]
\label{mm:thm:integral}
For every $P$ and positive $B,C$,
\begin{equation}\label{mm:eq:double-integral}
 \boxed{\displaystyle
 K_{P,a}(B,C)=\int_C^B\!\int_C^B
              \frac{R_{P,a}(u)-R_{P,a}(v)}{v-u}\,du\,dv.}
\end{equation}
The integrand at $u=v$ is $-R_{P,a}'(u)$.  In particular,
\begin{equation}\label{mm:eq:mixed-derivative}
 \partial_B\partial_CK_{P,a}(B,C)
       =-2\frac{R_{P,a}(B)-R_{P,a}(C)}{C-B}.
\end{equation}
The kernel and both of its first partial derivatives vanish on the
diagonal $B=C$.
\end{theorem}
\begin{proof}
At an ordinary lattice point,
\[
 \log^2\frac{x^2+B}{x^2+C}
   =\int_C^B\!\int_C^B\frac{du\,dv}{(x^2+u)(x^2+v)}.
\]
The rational kernel splits as
\[
 \frac1{(x^2+u)(x^2+v)}
   =\frac{(x^2+u)^{-1}-(x^2+v)^{-1}}{v-u}.
\]
For an arbitrary polynomial weight, perform the same finite
subtractions as in \eqref{mm:eq:K}.  Uniformity on the compact mass
interval permits both integrals to pass through the ordinary remainder;
the finitely many polynomial terms integrate directly.  The resulting
continued rational sum is the divided difference of
\eqref{mm:eq:RP}, proving \eqref{mm:eq:double-integral}.
The diagonal value is removable by differentiation.  Differentiating
the original logarithmic kernel and its finite subtractions gives
\eqref{mm:eq:mixed-derivative}.  The diagonal normalizations follow
from the square or from the double integral.
\end{proof}

\subsection{Every coincident-mass coefficient is finite}

Write $Z_{P}(m;B,a)$ for the continued value of
$\sum_nP(x_n^2)(x_n^2+B)^{-m}$.  It is regular at every positive
integer $m$, since possible spectral poles for an even polynomial
weight occur at half-integers.

\begin{theorem}[The full confluent expansion]
\label{mm:thm:confluence}
For $|h|<a^2+B$,
\begin{equation}\label{mm:eq:confluence}
 \boxed{\displaystyle
 K_{P,a}(B,B+h)=2\sum_{m\ge2}
       \frac{(-1)^mH_{m-1}}m\,h^m Z_{P}(m;B,a).}
\end{equation}
Every coefficient has the finite evaluation
\begin{equation}\label{mm:eq:coefficient-resolvent}
 Z_{P}(m;B,a)=\frac{(-1)^{m-1}}{(m-1)!}
                            \partial_B^{m-1}R_{P,a}(B).
\end{equation}
Consequently no new unspecified constants are introduced in this
coincident-mass expansion: the coefficients belong to the finite
algebra generated by polygamma values at $a\pm i\sqrt B$,
Bernoulli-polynomial values, and rational functions of $B$.
\end{theorem}
\begin{proof}
Use the ordinary Taylor identity, with the logarithm that vanishes at zero,
\[
 \log^2(1+z)=2\sum_{m\ge2}\frac{(-1)^mH_{m-1}}m z^m,
 \qquad |z|<1.
\]
The nearest possible zero of $x_n^2+B+h$ occurs at
$h=-(a^2+B)$.  On every smaller disk the finite-subtracted logarithmic
sum is holomorphic and the same subtractions justify coefficient
extraction.  More explicitly, separate the finitely many indices
$m\le D$, which use continued values, from $m\ge D+1$,
whose weighted sums converge ordinarily and normally on the smaller
disk.  A coefficient is precisely the continued weighted
resolvent $Z_P(m;B,a)$, which proves \eqref{mm:eq:confluence}.
Finally, differentiate the identity
$\partial_B Z_P(s;B,a)=-sZ_P(s+1;B,a)$, first in its convergence
region and then by continuation, starting at $s=1$.  No integer-order
pole is crossed.  This gives \eqref{mm:eq:coefficient-resolvent} and
the asserted finite special-function description.
\end{proof}

Equivalently, polynomial Taylor expansion about $-B$ gives
\[
 Z_P(m;B,a)=\sum_{\ell=0}^{D}\frac{P^{(\ell)}(-B)}{\ell!}
                          Z_1(m-\ell;B,a).
\]
For $j\ge0$ the required nonpositive-order values are simply
\[
 Z_1(-j;B,a)=\sum_{\ell=0}^{j}\binom j\ell
                              B^{j-\ell}\zeta(-2\ell,a),
\]
and the positive orders are derivatives of \eqref{mm:eq:R}.
This second formula is convenient for exact symbolic implementation.

For instance, when $P=1$,
\[
 K_{1,a}(B,B+h)
 =h^2Z_1(2;B,a)-h^3Z_1(3;B,a)
       +\frac{11}{12}h^4Z_1(4;B,a)+O(h^5).
\]
At $a=1/2$, $B=b^2$, the first coefficient is elementary:
\begin{equation}\label{mm:eq:half-explicit}
 Z_1(2;b^2,1/2)
 =\frac{\pi\tanh(\pi b)}{4b^3}
                     -\frac{\pi^2\operatorname{sech}^2(\pi b)}{4b^2}.
\end{equation}
It follows by differentiating
$R_{1/2}(b^2)=\pi\tanh(\pi b)/(2b)$.
The corresponding bilateral results follow by adding the formulas
at $a$ and $1-a$, for $0<a<1$.  In that sum all nonpositive even
Hurwitz values cancel by Bernoulli reflection, leaving an even shorter
weighted resolvent.  The higher spectral derivatives at a generic
noncoincident mass remain described by the exact logarithmic kernels;
no unproved finite polylogarithm reduction is inferred from the finite
confluent coefficients.

